%! TEX program = xelatex
%! TEX root = ../main.tex
\documentclass[../main]{subfiles}
\begin{document}
\section{Acuerdo aulico}
\subsection*{Normas para la presentación de la carpeta de trabajos
prácticos, teoría y acuerdos.}
\begin{enumerate}
  \item La portada debe ser la propuesta por el docente, en caso que se
    encuentra la portada de la escuela se utilizará la misma, o elaborar
    una que contenga la misma información, el índice debe estar completo,
    con la numeración completa, en la misma debe estar el programa de la
    materia y la presente guía de trabajos prácticos y pautas que deben
    estar firmadas por el alumno.
  \item Rótulos completos, realizados en regla, con todos los datos
    correspondientes y el formato propuesto por el docente, las hojas deben
    estar numeradas, letra técnica o imprenta prolija según lo pactado
    por el docente.
  \item La carpeta debe estar prolija, se tendrá en cuenta la ortografía,
    las hojas deben ser la misma sen cada trabajo práctico,
    al cambiar de trabajo pueden cambiar las hojas, se debe utilizar
    el mismo color de lapicera, los dibujos deben ser realizados en regla.
  \item Carpeta completa y entrega en termino.
  \item \textbf{Actitudinal:} El alumno deberá tener un vocabulario
    adecuado en el aula, respeto hacia el docente, hacia sus
    compañeros y las normas de convivencia.
  \item Trabajo en clase, trabajo en grupo, en pizarrón: Cada
    participación en clases será registrada por el docente en una
    panilla para luego generar una nota de trabajo en clase
    correspondiente al período escolar, que dependerá de la
    participación, cantidad y calidad.
  \item Asistencia (80\%).
\end{enumerate}
\vspace{1cm}
\begin{center}
  Firma Alumno $\dots\dots\dots\dots\dots\dots\dots\dots\dots\dots$
  \vspace{1cm}
  Firma Padre o Tutor $\dots\dots\dots\dots\dots\dots\dots\dots\dots\dots$
\end{center}
\end{document}
